\documentclass[12pt,a4paper]{article}
\usepackage[utf8]{inputenc}
\usepackage[greek]{babel} % for greek
%\usepackage[T1]{fontenc}
\usepackage{amsmath}
\usepackage[svgnames]{xcolor}
\usepackage{amsfonts}
\usepackage{amssymb}
\usepackage{graphicx}
\usepackage{hyperref}
\usepackage[left=1.50cm, right=1.50cm, top=3.00cm, bottom=3.00cm]{geometry}
%\author{Stratos Pateloudis}
\date{}
\title{Δικαιώματα Μαρίας-Νάσας}
\begin{document}
\maketitle
Τα συνολικά δέντρα είναι 185 (134 στο πάνω χωράφι και 51 στο κάτω, μετρημένα). Τα 51 δέντρα αντιστοιχούν σε ποσοστό $51/185=0.275$ των συνολικών δέντρων. Τα κιλά του λαδιού φέτος (η οποία ήταν η καλύτερη χρονιά μέχρι τώρα απο άποψη παραγωγής) ήταν 852.7 κιλά. 

\section{Σενάριο 1: ποσοστιαίος μοιρασμός}
Σε αυτό το σενάριο υποθέτουμε ότι η Νάσα και η Μαρία θα πάρουνε ποσοστιαία ότι τους αναλογεί σε σχέση με το χωράφι τους. Δηλαδή από την συνολική παραγωγή, η παραγωγή του χωραφιού τους είναι 51/185 επί της συνολικής παραγωγής. Αυτό βέβαια προϋποθέτει ότι και τα δύο χωράφια αποδίδουν το ίδιο. Ας θεωρήσουμε αυτό το σενάριο. 

Το λάδι που αντιστοιχεί στο κάτω χωράφι υπολογίζεται ως
\begin{equation}\nonumber
51/185\times 852.7=235.069~kg.
\end{equation} 
Αφαιρώντας τα δικαιώματα του ελαιοτριβείου φέτος (20\%) έχουμε υπόλοιπο
\begin{equation}\nonumber
235.069-0.2\times235.069=188.055~kg.
\end{equation} 

Αυτή είναι η τελική παραγωγή του κάτω χωραφιού που <<μπαίνει στα δοχεία>>, υποθέτωντας πάντα ισανάλογη παραγωγή και των δύο χωραφιών. Αν προτείνουμε το 1/3 να πηγαίνει σε Νάσα και Μαρία συνολικά, δικαιούνται από 
\begin{equation}\nonumber
188.055\times\frac{1}{3}\times\frac{1}{2}= 31.34~kg 
\end{equation} 
η κάθε μία που αντιστοιχεί σε 2.1 δεκαπεντάκιλους τενεκέδες.

\section{Σενάριο 2: με βάση την παραγωγή του χωραφιού.}
Υποθέτωντας οτι τα δύο χωράφια δεν δίνουν αναλογικά την ίδια παραγωγή έχει περισσότερο νόημα να προτείνουμε το 1/3 της παραγωγής του χωραφιού τους. Υποθέτωντας ότι το κάτω χωράφι έδωσε 1000 κιλά φέτος απο τα 4348 κιλά συνολικά, χρησιμοποιούμε την φετινή απόδοση για να βρούμε το λάδι. Δηλαδή στο κάτω χωράφι έχουμε 
 \begin{equation}\nonumber
 1000\times 0.196=196~kg. 
 \end{equation} 
 Ακολουθώντας την ίδια διαδικασία αφαιρούμε τα δικαιώματα του ελαιοτριβείου και έχουμε 
 \begin{equation}\nonumber
 196-0.2\times196=156.8~kg.
 \end{equation} 
Για την προσφορά του 1/3 της παραγωγής του χωραφιού έχουμε πάλι το τελικό ποσό
\begin{equation}\nonumber
156.8\times\frac{1}{3}\times\frac{1}{2}= 26.13~kg 
\end{equation} 
στην κάθε μια αντιστοιχώντας σε 1.74 δεκαπεντάκιλους τενεκέδες. 
\end{document}
